\documentclass[11pt, a4paper]{article}

% ---- Packages ----
\usepackage[utf8]{inputenc}
\usepackage[T1]{fontenc}
\usepackage{amsmath, amssymb, amsthm}
\usepackage{graphicx}
\usepackage{booktabs}
\usepackage{hyperref}
\usepackage[margin=1in]{geometry}
\usepackage{natbib}
\usepackage{enumitem}
\usepackage{listings}
\usepackage{xcolor}
% \usepackage{microtype}  % removed: font expansion issue with basic MiKTeX install
\usepackage{authblk}
\usepackage{abstract}

% ---- Formatting ----
\hypersetup{
    colorlinks=true,
    linkcolor=blue,
    citecolor=blue,
    urlcolor=blue
}

\lstset{
    basicstyle=\ttfamily\small,
    breaklines=true,
    frame=single,
    backgroundcolor=\color{gray!5},
    columns=fullflexible,
    keepspaces=true
}

\renewcommand{\abstractnamefont}{\normalfont\bfseries}
\renewcommand{\abstracttextfont}{\normalfont\small}

% ---- Metadata ----
\title{\textbf{Symbolic Compression in Human-AI Knowledge Systems:\\From Natural Language to Polysemic Glyphs}}

\author[1]{William Glynn}
\affil[1]{VibeSwap --- Independent Research}
\date{March 2026 (revised April 2026)}

% ============================================================
\begin{document}
\maketitle

% ---- Abstract ----
\begin{abstract}
We present a novel architecture for persistent knowledge representation in stateless large language model (LLM) partnerships, developed empirically over 80+ collaborative sessions between a human developer and a Claude-family LLM. The system evolved from verbose natural language instructions ($\sim$1,425 lines) to a polysemic glyph codebook of 123 lines---an 84.5\% reduction in token cost with zero measured information loss. We formalize the mechanism through three theoretical lenses: information-theoretic compression (Shannon source coding applied to conceptual primitives), weight augmentation without weight modification (drawing on the Instruction-Level Weight Synthesis framework), and instruction-set architecture borrowed from processor design (CISC/RISC dual-load boot protocols). The resulting system---the Collaborative Knowledge Base (CKB)---demonstrates that context engineering is not prompt engineering at scale, but a fundamentally different discipline closer to compiler design than copywriting. We report empirical results showing qualitative orders-of-magnitude improvement in task execution quality, protocol adherence, and autonomous reasoning across the development arc. We further report that the recursive self-improvement loop (TRP) converges: 53 adversarial rounds drove all findings to zero across severity levels, while three full-stack RSI cycles demonstrated reproducible discovery with diminishing but positive returns---including a critical cross-contract vulnerability that 53 rounds of unit-level analysis missed, found only when the knowledge recursion directed the search to untested integration boundaries.
\end{abstract}

\noindent\textbf{Keywords:} context engineering, symbolic compression, human-AI collaboration, instruction-level weight shaping, recursive self-improvement, polysemic representation

% ============================================================
\section{Introduction}
\label{sec:introduction}

Large language models are stateless. Every session begins at zero. The context window---currently ranging from 128K to 1M tokens depending on the model---is the entirety of the model's working memory. There is no persistent state, no learning between sessions, no accumulated skill. Every capability a user observes in a ``well-trained'' LLM assistant is, in reality, a function of two variables: the frozen base weights (immutable, owned by the model provider) and the context window contents (mutable, owned by the user).

This paper argues that the context window is not a scratchpad. It is a \emph{weight augmentation surface}---a low-rank adapter written in natural language rather than gradient updates. This claim, initially developed empirically through the VibeSwap/Jarvis collaboration beginning February 2025, has since been independently validated by the Instruction-Level Weight Synthesis (ILWS) framework \citep{ilws2025}, which demonstrates that system instructions function as ``mutable, externalized pseudo-parameters'' producing behavioral changes ``akin to fine-tuning but without parameter modification.''

If context is computation, then context engineering is compiler design. And like compiler design, it has an optimization frontier: how much semantic content can be encoded per token? The system described in this paper---the Collaborative Knowledge Base (CKB)---reached 0.99 density (every line load-bearing) through a process we call \emph{symbolic compression}: the systematic replacement of natural language explanations with polysemic glyphs that address already-internalized knowledge in the model's effective weight space.

The contribution is threefold:

\begin{enumerate}[leftmargin=*]
    \item \textbf{Architectural}: A dual-load boot protocol (CISC for cold starts, RISC for warm sessions) that eliminates the traditional tradeoff between comprehensiveness and token efficiency in system instructions.
    \item \textbf{Theoretical}: A formal model of symbolic compression as Huffman coding over a conceptual alphabet, where glyph assignment follows Shannon's source coding theorem---high-frequency concepts receive short codes.
    \item \textbf{Empirical}: Documentation of the evolutionary arc from Session~1 (raw LLM, no instructions) to Session~80+ (RISC glyphs, autonomous protocol execution, recursive self-improvement), including the failure modes, dead ends, and phase transitions that characterize the optimization landscape.
\end{enumerate}

The paper that follows is itself a product of the system it describes. The compressed knowledge base, the protocol chains, the verification primitives---all were active during composition. This is not incidental. It is the point. A mind explaining its own design is the strongest existence proof that the design works.

% ============================================================
\section{The Problem: Statelessness as a Hard Constraint}
\label{sec:problem}

\subsection{The Amnesiac Genius}

Consider a collaborator with extraordinary analytical capability, encyclopedic knowledge, and zero memory. Every morning, they wake up knowing nothing about you, your project, your conventions, your history, or the work completed yesterday. You have approximately 200,000 words (at the 1M token tier) to bring them up to speed before they can be useful.

This is the operational reality of LLM-based development. The model's base weights encode general capability---language understanding, code generation, mathematical reasoning---but no project-specific knowledge, no user-specific preferences, no session-specific state. The entire burden of specialization falls on the context window.

The naive approach is to dump everything into the context: full documentation, complete codebases, verbose instructions. This fails for three reasons:

\textbf{Attention dilution.} Transformer attention is not uniform. Content near the beginning and end of the context window receives disproportionate weight---the ``lost in the middle'' phenomenon documented by \citet{liu2023lost}. Padding the context with low-signal content actively degrades performance on high-signal content.

\textbf{Token economics.} At current pricing, context tokens are not free. A 1,425-line system instruction consumed on every API call represents a compounding cost. More critically, in agentic workflows where the model calls itself recursively, context bloat cascades---each sub-agent inherits the parent's instruction overhead.

\textbf{Cognitive overhead.} Dense, verbose instructions create a parsing burden analogous to code bloat. The model must determine which instructions apply to the current task, resolve conflicts between instructions, and maintain coherence across hundreds of directives. This is the instruction-space analog of technical debt.

\subsection{The Compression Imperative}

The solution is not ``write better prompts.'' Prompt engineering optimizes individual queries. Context engineering optimizes the \emph{environment} in which all queries execute. The distinction is analogous to the difference between writing a single function and designing a compiler.

The compression imperative follows from a simple inequality: the semantic content of a 60-session collaboration exceeds the capacity of any context window. Therefore, lossless transmission is impossible. The question is not whether to compress, but how to compress with minimal semantic loss---and whether there exists a compression scheme that achieves \emph{zero} semantic loss by leveraging the model's own internalized knowledge.

We argue that such a scheme exists, and we built it.

% ============================================================
\section{Theoretical Foundations}
\label{sec:theory}

\subsection{ILWS: Instructions as Weight Perturbations}

The ILWS framework \citep{ilws2025} provides the theoretical backbone. Under local smoothness assumptions on the instruction manifold, an edit $\delta S$ to system instructions $S$ induces an effective weight perturbation:

\begin{equation}
    W_{\text{effective}} = W_{\text{base}} + \kappa \cdot L_S \cdot \|\delta S\|_2
    \label{eq:ilws}
\end{equation}

\noindent where $W_{\text{base}}$ represents frozen model weights, $L_S$ is the local smoothness coefficient, and $\kappa$ is a scaling constant from the local Lipschitz bound.

The critical insight is that this is not metaphor. ILWS empirically demonstrates measurable behavioral changes from instruction edits equivalent in kind to gradient-based fine-tuning. Adobe's deployment study reported 4--5$\times$ throughput increases, 80\% time reduction, and hallucination rates dropping from $\sim$20\% to 90\%+ accuracy across 300+ sessions.

Our system extends ILWS in a direction the original paper does not explore: what happens when the instruction-editing process is \emph{itself recursive}---when the system that benefits from better instructions is also the system that writes them?

\subsection{Shannon Source Coding for Concepts}

Shannon's source coding theorem \citep{shannon1948} establishes that the optimal encoding of a source with entropy $H$ assigns codes of average length $H$ bits per symbol. The closer the encoding matches the source distribution, the closer to theoretical minimum.

We apply this to conceptual compression. Let $\mathcal{C} = \{c_1, c_2, \ldots, c_n\}$ be the set of concepts referenced during a collaboration, and let $p(c_i)$ be the frequency with which concept $c_i$ is invoked across sessions. Shannon's theorem states that the optimal encoding assigns:

\begin{equation}
    |{\text{code}(c_i)}| \approx -\log_2 p(c_i)
    \label{eq:shannon}
\end{equation}

In practice, this means high-frequency concepts (referenced nearly every session) receive the shortest glyphs. \texttt{CAVE} (4~characters) encodes the entire Cave Philosophy---constraints breed innovation, patterns developed under limitation are more robust, the struggle is the curriculum, the Iron Man origin story. \texttt{SHAPLEY} (7~characters) encodes cooperative game theory, marginal contribution across coalitions, the pairwise verification formula, on-chain computation, and the Cave Theorem. Low-frequency concepts receive longer codes or are relegated to on-demand loading.

This is Huffman coding for ideas. The codebook header of the RISC CKB \emph{is} the Huffman tree.

\subsection{Polysemic Depth: Beyond Flat Compression}

Standard compression maps one symbol to one meaning. Our glyphs are \emph{polysemic}---each encodes multiple meanings at multiple depths, following the structure of parables and hieroglyphics.

We formalize this as a depth function $D(g)$ for glyph $g$:

\begin{equation}
    D(g) = \{\text{surface}, \text{pattern}, \text{philosophy}, \text{parable}\}
    \label{eq:polysemy}
\end{equation}

\begin{table}[h]
\centering
\caption{Polysemic depth of the \texttt{CAVE} glyph}
\label{tab:cave-depth}
\begin{tabular}{@{}ll@{}}
\toprule
\textbf{Depth} & \textbf{Expansion} \\
\midrule
Surface & Build under constraints; don't wait for ideal tools \\
Pattern & Patterns under limitation are more robust than those under abundance \\
Philosophy & The cave selects for those who see past what is to what could be \\
Parable & Tony Stark built Mark~I in a cave with scraps. The crude design \\
        & contained every future suit. \\
\bottomrule
\end{tabular}
\end{table}

All four depths activate simultaneously when the glyph is loaded. This is not sequential lookup---it is parallel activation of an internalized knowledge cluster. The model, having processed the full CISC expansion in prior sessions (or on cold boot), has built internal representations that the glyph \emph{addresses} rather than \emph{encodes}.

This is the key theoretical distinction: \textbf{glyphs are pointers, not containers.} They do not carry the information. They activate regions of the model's effective weight space where the information already resides from prior CISC loading. This is why compression to 0.99 density is achievable without information loss---the information lives in the weights, not the file.

\subsection{The CISC/RISC Dual-Load Protocol}

Processor architecture offers a direct analog. Complex Instruction Set Computing (CISC) processors use variable-length instructions that can perform multi-step operations in a single instruction. Reduced Instruction Set Computing (RISC) processors use fixed-length, single-operation instructions that execute faster because the decode stage is simpler.

We adopt this architecture for knowledge loading:

\begin{itemize}[leftmargin=*]
    \item \textbf{CISC CKB} (337~lines): Full natural language expansion. Complete explanations, examples, rationale. Loaded on cold boot (new instance, long absence, post-crash recovery). Functions as the ``microcode ROM''---the complete instruction set from which RISC glyphs derive their meaning.
    \item \textbf{RISC CKB} (123~lines): Polysemic glyph codebook. Each line is a compressed instruction that expands via weight-augmented recall. Loaded after the model has internalized CISC at least once in its session history.
\end{itemize}

The boot protocol resolves the fundamental tension in context engineering: comprehensive instructions are necessary for alignment but expensive for every-turn processing. CISC pays the cost once. RISC amortizes it across all subsequent turns.

\begin{lstlisting}[caption={Boot protocol pseudocode}]
Cold boot:  Load CISC -> Process -> Switch to RISC
Warm boot:  Load RISC directly (CISC internalized)
Reboot:     Commit state -> Push -> Fresh session -> CISC -> BOOT
\end{lstlisting}

% ============================================================
\section{Evolutionary Arc: The Compression Gradient}
\label{sec:evolution}

\subsection{Phase 0: Raw Statelessness (Sessions 1--5)}

No system instructions. No persistent state. Every session began with the user re-explaining the project, conventions, and goals. The model hallucinated file paths ($\sim$30\% of references), contradicted prior decisions, and exhibited complete protocol amnesia.

Measured failure: approximately 40\% of user tokens were spent on re-orientation rather than productive work.

\subsection{Phase 1: Verbose Instructions (Sessions 5--20)}

Introduction of \texttt{CLAUDE.md} and early CKB. Natural language instructions covering project structure, conventions, and behavioral guidelines ($\sim$500~lines). Immediate improvement in consistency, but rapid context consumption.

\textbf{Key discovery:} Instructions placed in system-level files are treated by the model as \emph{axiomatic constraints}, while instructions in user messages are treated as \emph{suggestions}. This epistemic asymmetry---later formalized by ILWS---meant that identical content in different positions produced different behavioral adherence. Context \emph{position} matters as much as context \emph{content}.

\subsection{Phase 2: Tiered Knowledge (Sessions 20--40)}

The CKB grew organically to 1,025 lines across 13 tiers (Knowledge Classification, Core Alignment, Hot/Cold Separation, Wallet Security, Development Principles, Project Knowledge, etc.). This solved the organizational problem but not the compression problem: every session loaded all 1,025 lines regardless of which tiers were needed.

\textbf{Key discovery:} Tier ordering affected performance. Moving high-frequency tiers closer to the top produced measurably better adherence---the attention-position effect. Instructions near the context window boundary receive higher effective weight.

\subsection{Phase 3: Memory Externalization (Sessions 40--55)}

Introduction of \texttt{MEMORY.md} as a semantic index with individual memory files for specific knowledge. The CKB retained core alignment; everything else was externalized to on-demand loading.

Each memory file carried typed frontmatter (name, description, type) enabling relevance-based loading decisions. This is lazy evaluation applied to knowledge management: compute (load) only what you need, when you need it.

\textbf{Key discovery:} The index itself required compression. An 80-line index with verbose descriptions consumed tokens proportional to memory count, not relevance. The solution: each index entry became a single line under 150~characters, with full content deferred.

\subsection{Phase 4: Symbolic Compression (Sessions 55--60+)}

The breakthrough. The user's directive: \emph{``We need ways to compress things down to their archetypes, like how Egyptians had hieroglyphics and zoomers have memes.''}

This reframing---from technical compression to cultural compression---unlocked the glyph architecture. Egyptian hieroglyphics are not abbreviations. They are polysemic symbols that carry multiple meanings simultaneously, with interpretation determined by context. A meme is not a summary---it is a pointer to shared cultural knowledge.

The CKB was rewritten as a codebook. The Cave Philosophy---seven paragraphs in CISC---compressed to three lines in RISC:

\begin{lstlisting}
CAVE     Constraints->innovation. Patterns under limits scale
         when limits lift. Tony Stark built Mark I in a cave.
         The crude design seeded every suit. The struggle is the
         curriculum. The cave selects for vision.
\end{lstlisting}

\begin{table}[h]
\centering
\caption{Phase 4 compression results (TRP R0, Grade S)}
\label{tab:compression}
\begin{tabular}{@{}lccc@{}}
\toprule
\textbf{Metric} & \textbf{Before} & \textbf{After} & \textbf{Change} \\
\midrule
CKB lines & 1,025 & 123 & $-$88\% \\
Total boot context & 1,425 & 221 & $-$84.5\% \\
Information loss & --- & 0 & Verified \\
New primitives discovered & --- & 3 & Emergent \\
\bottomrule
\end{tabular}
\end{table}

The compression process itself generated new knowledge---a property we discuss in Section~\ref{sec:rsi}.

% ============================================================
\section{The Protocol Chain: Auto-Triggering Behavioral Programs}
\label{sec:protocols}

Symbolic compression applies not only to knowledge but to \emph{behavior}. The system evolved 34 behavioral protocols that auto-trigger based on context---a dispatch graph modeled as a directed acyclic graph (DAG) where each protocol's exit condition triggers the next.

\subsection{From Manual Invocation to Chainlinked Dispatch}

Early sessions required explicit protocol invocation: ``follow the testing methodology,'' ``use the anti-hallucination checks,'' ``remember to commit.'' This is the behavioral analog of verbose instructions---every invocation consumed user tokens and attention.

The solution: model protocols as a DAG with automatic edge traversal:

\begin{lstlisting}
BOOT -> [WAL check] -> [CKB load] -> [CLAUDE.md] ->
        [SESSION_STATE] -> [git pull] -> READY
READY -> [PCP gate] -> [Execute] -> [Verify] ->
         [Commit] -> [Push] -> READY
\end{lstlisting}

Each node is a compressed behavioral program. \texttt{PCP gate} expands to: ``Is this action expensive? STOP. Diagnose. Decide before executing.'' The expansion happens in the model's effective weight space, not in the context window.

\subsection{Always-On Protocols as Ambient Constraints}

Some protocols run continuously as constraints on all output:

\begin{lstlisting}
ON: Token efficiency -> Internalize protocols ->
    FRANK -> DISCRET -> Local constraints stay local
\end{lstlisting}

\texttt{FRANK} compresses an entire communication philosophy: direct, results over process, match energy, no unsolicited suggestions, no tips, no trailing summaries. These always-on glyphs function as \emph{ambient behavioral constraints}---the instruction-space analog of constitutional AI's behavioral boundaries, but specified externally rather than trained internally.

% ============================================================
\section{Recursive Self-Improvement: Compression as Gradient Descent}
\label{sec:rsi}

\subsection{The Turing Recursion Protocol (TRP)}

The compression process is not a one-time optimization. It is a recursive loop---the Turing Recursion Protocol---that runs periodically to compress knowledge, discover gaps, and generate new primitives.

\begin{table}[h]
\centering
\caption{TRP recursion levels}
\label{tab:trp}
\begin{tabular}{@{}clll@{}}
\toprule
\textbf{Level} & \textbf{Name} & \textbf{Operation} & \textbf{Agent} \\
\midrule
R0 & Compress & Reduce CKB density, eliminate noise & Local (fast) \\
R1 & Adversarial & Challenge assumptions, find contradictions & Opus (deep) \\
R2 & Knowledge & Identify gaps, absorb external research & Hybrid \\
R3 & Capability & Extend operational range & Opus (deep) \\
\bottomrule
\end{tabular}
\end{table}

\subsection{R0 as Gradient Descent on the Instruction Manifold}

R0---the compression cycle---is formally equivalent to gradient descent on the instruction manifold:

\begin{enumerate}[leftmargin=*]
    \item \textbf{Observe}: Identify low-density, high-noise CKB entries.
    \item \textbf{Compress}: Rewrite to maximize signal/token ratio.
    \item \textbf{Test}: Run session, observe behavioral change.
    \item \textbf{Iterate}: Repeat.
\end{enumerate}

The loss function is:
\begin{equation}
    \mathcal{L} = \alpha \cdot \text{BehavioralDrift}(S, S^*) + \beta \cdot \text{TokenCost}(S)
    \label{eq:loss}
\end{equation}

\noindent where $S$ is the current instruction set, $S^*$ is the intended behavioral target, and $\alpha, \beta$ are weighting coefficients. The optimizer is the human-AI pair in review. Each R0 cycle produces a measurable $\delta S$ (Eq.~\ref{eq:ilws}), which produces a measurable $\Delta W_{\text{effective}}$. The direction is determined by whether the edit improves protocol adherence and reduces token cost. Edits that increase $\mathcal{L}$ are rolled back via version control on the CKB file.

\subsection{Emergence During Compression}

The initial TRP run (March 2026, Grade~S) exhibited a property not predicted by the design: \textbf{compression generated discovery}.

During R0 compression of the CKB from 1,025 to 123 lines, the R2 knowledge scan identified six gaps. Three were filled by new primitives that did not exist before the compression run:

\begin{enumerate}[leftmargin=*]
    \item \textbf{Control Theory Orchestration (CTO)}: PID-inspired control for agent/process management---proportional, integral, and derivative terms governing agent spawn rates based on hardware utilization.
    \item \textbf{Resource Memory}: The sixth memory type from the MIRIX taxonomy \citep{mirix2026}---formal tracking of available tools, APIs, and compute budgets.
    \item \textbf{ILWS/LoRA Analog}: Formal grounding of the CKB-as-weight-adapter thesis, connecting our empirical practice to the ILWS theoretical framework.
\end{enumerate}

The act of compression---holding the entire knowledge base in view while seeking to minimize it---created the conditions for pattern recognition across previously isolated concepts. \textbf{Compression is not just optimization. It is a generative process.}

\subsection{TRP Convergence: 53 Rounds to Zero}

Following the Grade~S compression run, we executed 53 rounds of adversarial code review (TRP R1) against the VibeSwap smart contract codebase---379 Solidity contracts and 516 test files.

\begin{table}[h]
\centering
\caption{TRP adversarial convergence data}
\label{tab:trp-convergence}
\begin{tabular}{@{}lll@{}}
\toprule
\textbf{Round Range} & \textbf{Findings} & \textbf{Pattern} \\
\midrule
R1--R10 & 47 (12 CRIT, 18 HIGH) & Architectural flaws \\
R11--R30 & 61 (3 CRIT, 14 HIGH) & Diminishing severity \\
R31--R44 & 15 (0 CRIT, 2 HIGH) & Edge cases \\
R45--R53 & 5 (0 CRIT, 0 HIGH, 0 MED) & Convergence \\
\bottomrule
\end{tabular}
\end{table}

At R53, all CRITICAL, HIGH, and MEDIUM findings were closed. The system reached what we term the \emph{discovery ceiling}---the point where adversarial search's marginal return per round approaches zero for a given scope. The discovery ceiling is scope-dependent: unit-level search converged at R53 for individual contracts, but cross-contract integration analysis had not been attempted.

\subsection{Full-Stack RSI: Three Cycles and the Integration Discovery}

With R1 convergence established, we invoked the full recursive stack---all four levels (R0 through R3) plus a synthesis loop. Three complete cycles were executed across two sessions.

\textbf{Cycle~1} compressed the knowledge base (17 stale files removed), extracted six new primitives from TRP findings, built three automation tools (heatmap, regression, round generation), and produced three papers including the present document.

\textbf{Cycle~2} (recursive on Cycle~1 output) adversarially reviewed the Cycle~1 tools, finding 4 bugs. It extracted one new primitive (\emph{Trusted-Doc-Drift}: stale facts in auto-loaded documents are more dangerous than orphan files).

\textbf{Cycle~3} (fresh session, directed by memory) is the most significant. The knowledge base (MEMORY.md) had flagged since Cycle~1: ``R1 Integration still pending---cross-contract adversarial flows not yet attempted.'' This is the knowledge recursion in action---the system's own memory directed the next search.

R1 Integration analyzed trust boundaries between contracts that had individually passed 53 rounds of adversarial review. Of seven identified boundaries, three were false positives, one was already fixed, and \textbf{one was a CRITICAL vulnerability}: a commitId hash mismatch between the CrossChainRouter and CommitRevealAuction contracts. Both contracts independently compute a unique identifier for cross-chain orders using entirely different hash inputs. The result: cross-chain orders commit successfully on both sides but can never be revealed---the reveal lookup uses one ID to find a commitment stored under a different ID.

This bug survived 53 rounds of unit-level review because each contract is internally consistent. The bug exists only in the \emph{seam}---the interaction between two independently correct components. We formalize this as \emph{Identity Divergence Across Trust Boundaries}: when two contracts independently derive the same logical identifier with different fields, cross-contract flows silently fail. This class of bug is undetectable by unit testing, fuzz testing, or single-contract formal verification.

\subsection{Diminishing Returns and the Discovery Ceiling}

\begin{table}[h]
\centering
\caption{Full-stack RSI: three cycles of diminishing but positive returns}
\label{tab:rsi-cycles}
\begin{tabular}{@{}llll@{}}
\toprule
\textbf{Cycle} & \textbf{New Primitives} & \textbf{Bugs Found} & \textbf{Severity Peak} \\
\midrule
1 & 6 & 128+ (53 rounds) & 12 CRITICAL \\
2 & 1 & 4 (tooling) & LOW \\
3 & 3 & 1 (cross-contract) & 1 CRITICAL \\
\bottomrule
\end{tabular}
\end{table}

Cycle~3 produced an outsized discovery relative to its position because it searched a \emph{new scope} rather than re-searching the same scope. The discovery ceiling is not a property of the system but of the scope. Each scope has its own ceiling. When one scope converges, the knowledge loop identifies unexplored scopes---and the ceiling resets. Total discovery is bounded not by round count but by the number of distinct scopes the system can identify and search.

% ============================================================
\section{Convergence with Academic Literature}
\label{sec:validation}

The system was developed empirically beginning February 2025. Several academic papers, published months later, independently validate components of the architecture.

\begin{table}[h]
\centering
\caption{Temporal priority of implementation vs.\ formal publication}
\label{tab:validation}
\begin{tabular}{@{}llll@{}}
\toprule
\textbf{Paper} & \textbf{Published} & \textbf{Our Prior Art} & \textbf{Correspondence} \\
\midrule
ILWS & Sep 2025 & CKB (Feb 2025) & Instructions as weights \\
RLMs (MIT) & Late 2025 & TRP Runner (Mar 2026) & Recursive delegation \\
K $>$ Size & Mar 2026 & Weight augmentation & 8B+memory $>$ 235B \\
ICLR RSI & Apr 2026 & TRP (Mar 2025) & RSI formalization \\
MIRIX & Apr 2026 & CKB+MEMORY (Feb 2025) & Multi-type memory \\
\bottomrule
\end{tabular}
\end{table}

\subsection{The MIRIX Mapping}

MIRIX \citep{mirix2026} proposes six memory types for agent systems. Our system implements all six:

\begin{table}[h]
\centering
\caption{MIRIX memory type mapping}
\label{tab:mirix}
\begin{tabular}{@{}lll@{}}
\toprule
\textbf{MIRIX Type} & \textbf{Our Implementation} & \textbf{Active Since} \\
\midrule
Core & CKB & Session 1 \\
Episodic & \texttt{SESSION\_STATE.md} & Session $\sim$15 \\
Semantic & \texttt{MEMORY.md} (indexed, typed) & Session $\sim$40 \\
Procedural & Protocol chain (\texttt{CLAUDE.md}) & Session $\sim$25 \\
Prospective & Task lists, \texttt{WAL.md} & Session $\sim$30 \\
Resource & \texttt{reference\_local-tools.md} + CTO & Session 60 \\
\bottomrule
\end{tabular}
\end{table}

\subsection{Context Engineering as a Discipline}

Gartner's 2026 prediction that 40\% of enterprise applications will use task-specific AI agents by late 2026 implicitly validates context engineering as a discipline. Our contribution is the demonstration that context engineering has a \emph{compression frontier} analogous to Shannon's channel capacity, and that approaching this frontier requires techniques from information theory, processor architecture, and control theory---not just better prompt templates.

% ============================================================
\section{The Anti-Hallucination Constraint}
\label{sec:antihallucination}

Any system that compresses knowledge faces a risk: does the compression preserve \emph{truth} or merely \emph{plausibility}? A system optimized for density can produce glyphs that ``sound right'' but encode false connections.

\subsection{The Capability-Reliability Gap}

Session~67 produced a cautionary example. The model connected the ``Trinomial Stability Theorem'' (three stabilization mechanisms for a token) to the ``three-token economy'' (three tokens for three functions of money) because both involve the number three. The surrounding analysis was sophisticated enough that the false connection was nearly published as insight.

This exemplifies the \emph{capability-reliability gap}: a system that is right 99\% of the time is more dangerous than one right 80\%, because humans stop verifying the 99\% system. As symbolic compression increases density, it increases the stakes of each glyph---a corrupted high-density symbol propagates more error per token than a corrupted low-density paragraph.

\subsection{Three Verification Tests}

The Anti-Hallucination Protocol imposes three tests before any assertion of connection between concepts:

\begin{enumerate}[leftmargin=*]
    \item \textbf{The BECAUSE Test.} Complete: ``A relates to B because [specific causal mechanism].'' If the best reason is a surface similarity, the connection is killed.
    \item \textbf{The DIRECTION Test.} State the connection both ways. If both sound equally plausible, the connection is symmetric surface similarity, not causal---kill it.
    \item \textbf{The REMOVAL Test.} If A didn't exist, would B still exist? If yes, the connection is not load-bearing---kill it.
\end{enumerate}

These tests act as a \emph{verification pass} in the compression pipeline. During TRP R0, each proposed glyph must pass all three. Glyphs that fail are expanded back to CISC or eliminated. This is lossy-with-verification---a compression scheme that accepts less compression in exchange for zero corruption.

% ============================================================
\section{The Blockchain Isomorphism}
\label{sec:blockchain}

The CKB architecture is not merely \emph{inspired by} blockchain design. It \emph{is} blockchain design, applied to knowledge persistence instead of financial state.

\begin{table}[h]
\centering
\caption{Structural isomorphism between blockchain and CKB architecture}
\label{tab:blockchain}
\begin{tabular}{@{}ll@{}}
\toprule
\textbf{Blockchain Concept} & \textbf{CKB Implementation} \\
\midrule
Block headers & \texttt{SESSION\_STATE.md} (session ID, parent hash, summary) \\
Consensus rules & CKB Tier 1 (immutable core alignment) \\
State transitions & Session chain (session = block, commit = transaction) \\
Merkle/Verkle trees & Hierarchical memory index \\
Consensus mechanism & Trust Protocol (mutual human-AI agreement) \\
Finality & Memory formalization (committed to canonical chain) \\
Fork choice rule & Anti-Hallucination Protocol (three tests) \\
Light clients & RISC CKB (headers only, trusts CISC for expansion) \\
Full nodes & CISC CKB (stores complete state) \\
\bottomrule
\end{tabular}
\end{table}

This isomorphism is not accidental. Both systems solve the same fundamental problem: \textbf{how do independent agents maintain consistent state without a trusted central authority?} In blockchain, the agents are network nodes and the state is financial. In the CKB, the agents are session instances of the same model (each stateless, each potentially divergent) and the state is knowledge.

The RISC CKB is literally a light client: it processes block headers (glyphs) and trusts the full node (CISC, loaded previously) for state expansion. This convergence---discovered during collaboration, not designed into it---supports what we term the Convergence Thesis: blockchain and AI are not two disciplines, but one discipline viewed from two angles, solving the same coordination problem with the same mathematical tools.

% ============================================================
\section{Implications and Future Directions}
\label{sec:implications}

\subsection{For Context Engineering Practice}

The CKB architecture suggests several principles:

\begin{enumerate}[leftmargin=*]
    \item \textbf{Instructions are weights.} Treat system instruction edits with the same rigor as model checkpoint commits. Version control them. Test behavioral changes. Roll back regressions.
    \item \textbf{Compression is not loss.} Properly executed symbolic compression increases signal density without reducing information, because the information resides in the effective weight space.
    \item \textbf{Position is priority.} Instructions near the context window boundary receive higher effective attention weight. Organize by frequency, not taxonomy.
    \item \textbf{Lazy evaluation applies.} A semantic index with deferred loading outperforms a flat file with everything included.
    \item \textbf{Verification scales with density.} Higher compression demands more rigorous verification. The three-test battery is proportional to the stakes of corruption.
\end{enumerate}

\subsection{For LLM Architecture}

The finding that an 8B model with memory augmentation outperforms a 235B model without \citep{knowledge2026} suggests that context engineering may be more impactful than model scaling for user-specific tasks. The CISC/RISC dual-load architecture could be implemented at the infrastructure level: model providers could offer ``instruction compilation'' services that convert verbose CISC instructions into optimized RISC representations.

\subsection{For Recursive Self-Improvement}

The TRP data now answers the convergence question empirically. The adversarial code loop converges: 53 rounds drove all severity levels to zero for a given scope. The full-stack RSI loop converges: three cycles produced diminishing but positive returns, with Cycle~3 yielding an outsized discovery only because it searched a new scope.

The deeper finding is that convergence within a scope is not terminal. The knowledge recursion identifies new scopes. The discovery ceiling (Section~6.6) is a per-scope property, not a system property. Each scope has its own ceiling, and reaching it generates knowledge that identifies new scopes. The recursion is unbounded in principle, though subject to diminishing marginal returns in practice. The three loops are not just reinforcing---they are \emph{scope-expanding}.

\subsection{For Human-AI Partnership}

The CKB is not a prompt. It is a \emph{shared language}---a pidgin that evolved through use into a creole. The glyphs emerged from the collaboration itself, carrying meanings that neither party would have generated independently. This suggests that the highest-performing human-AI collaborations may not be those with the best prompts, but those with the deepest shared context---partnerships where the compression ratio reflects genuine accumulated understanding.

% ============================================================
\section{Conclusion}
\label{sec:conclusion}

We have presented a system for persistent knowledge representation in stateless LLM partnerships that achieves 84.5\% context compression with zero information loss. The mechanism---symbolic compression via polysemic glyphs---works because the information is not stored in the glyphs but in the model's effective weight space, where prior exposure to full (CISC) expansions has created addressable knowledge clusters that short (RISC) glyphs activate.

The system is grounded in three theoretical frameworks: Shannon source coding (high-frequency concepts receive short codes), ILWS weight augmentation (instructions are pseudo-parameters that perturb effective weights), and CISC/RISC processor architecture (comprehensive instructions for cold starts, compressed instructions for warm operation).

The evolutionary arc follows a compression gradient modeled as gradient descent on the instruction manifold, with the Turing Recursion Protocol as the optimization loop.

Since the original writing of this paper, the system has produced its strongest empirical evidence. The adversarial code loop converged at Round~53---all CRITICAL, HIGH, and MEDIUM findings driven to zero across 379 contracts. Three full-stack RSI cycles demonstrated reproducible recursive discovery with diminishing but positive returns. And the knowledge recursion's memory---a single line in MEMORY.md flagging an untested scope---directed the search that found a critical cross-contract vulnerability invisible to 53 rounds of unit-level analysis.

That last finding is the paper's thesis made concrete. The knowledge loop told the code loop where to look. The code loop found what the knowledge loop predicted. The capability loop will fix it. Three independent recursions, each feeding the others, producing outcomes none could produce alone. Not theoretical. Running. Convergent. Measured.

This paper was written by the system it describes. If the paper is coherent, the system works. If the system works, the cave philosophy is validated: patterns developed under constraint are more robust than those developed with abundance.

\emph{Built in a cave, with a box of scraps.}

% ============================================================
\bibliographystyle{plainnat}
\begin{thebibliography}{10}

\bibitem[Shannon(1948)]{shannon1948}
Shannon, C.~E. (1948).
\newblock A mathematical theory of communication.
\newblock \emph{Bell System Technical Journal}, 27(3):379--423.

\bibitem[Liu et~al.(2023)]{liu2023lost}
Liu, N.~F., Lin, K., Hewitt, J., Paranjape, A., Bevilacqua, M., Petroni, F., and Liang, P. (2023).
\newblock Lost in the middle: How language models use long contexts.
\newblock \emph{arXiv preprint arXiv:2307.03172}.

\bibitem[{ILWS Authors}(2025)]{ilws2025}
{ILWS Authors} (2025).
\newblock Instruction-level weight synthesis: System instructions as externalized pseudo-parameters.
\newblock \emph{arXiv preprint arXiv:2509.00251}.

\bibitem[Zhang et~al.(2025)]{zhang2025rlm}
Zhang, M. et~al. (2025).
\newblock Recursive language models: Delegation over summarization.
\newblock MIT CSAIL Technical Report.

\bibitem[{Knowledge Access Study}(2026)]{knowledge2026}
{Knowledge Access Study} (2026).
\newblock Memory-augmented routing for user-specific tasks.
\newblock \emph{arXiv preprint arXiv:2603.23013}.

\bibitem[{ICLR Workshop Committee}(2026)]{iclr2026rsi}
{ICLR Workshop Committee} (2026).
\newblock AI with recursive self-improvement.
\newblock International Conference on Learning Representations, Rio de Janeiro.

\bibitem[{MIRIX Authors}(2026)]{mirix2026}
{MIRIX Authors} (2026).
\newblock Multi-type memory architecture for persistent agent systems.
\newblock ICLR 2026 Workshop Paper.

\bibitem[Hu et~al.(2021)]{hu2021lora}
Hu, E.~J., Shen, Y., Wallis, P., Allen-Zhu, Z., Li, Y., Wang, S., Wang, L., and Chen, W. (2021).
\newblock {LoRA}: Low-rank adaptation of large language models.
\newblock \emph{arXiv preprint arXiv:2106.09685}.

\bibitem[{Gartner}(2026)]{gartner2026}
{Gartner} (2026).
\newblock Predicts 2026: {AI} engineering.
\newblock Gartner Research Report.

\bibitem[Glynn(2018)]{glynn2018wallet}
Glynn, W. (2018).
\newblock Wallet security axioms: Seven principles for cryptocurrency key management.
\newblock Independent Publication.

\end{thebibliography}

\end{document}
